% chebpy.tex --- A description of the ChebPy library, its architecture and
% the numerical algorithms it implements.
%
% Compiled by .github/workflows/rhiza_paper.yml via
%     latexmk -pdf -bibtex -interaction=nonstopmode chebpy.tex
% or locally with `make paper`.

\documentclass[11pt,a4paper]{article}

\usepackage[T1]{fontenc}
\usepackage[utf8]{inputenc}
\usepackage{lmodern}
\usepackage{microtype}
\usepackage[margin=2.7cm]{geometry}
\usepackage{amsmath}
\usepackage{amssymb}
\usepackage{amsthm}
\usepackage{booktabs}
\usepackage{array}
\usepackage{listings}
\usepackage{xcolor}
\usepackage[hidelinks]{hyperref}

\newtheorem{theorem}{Theorem}
\newtheorem{proposition}[theorem]{Proposition}

\definecolor{codebg}{gray}{0.96}
\definecolor{codekw}{rgb}{0.13,0.29,0.53}
\definecolor{codestr}{rgb}{0.16,0.44,0.16}
\definecolor{codecom}{gray}{0.45}

\lstdefinestyle{py}{
  language=Python,
  basicstyle=\ttfamily\small,
  keywordstyle=\color{codekw}\bfseries,
  stringstyle=\color{codestr},
  commentstyle=\color{codecom}\itshape,
  backgroundcolor=\color{codebg},
  frame=single,
  rulecolor=\color{codebg},
  framesep=5pt,
  showstringspaces=false,
  columns=fullflexible,
  breaklines=true,
  xleftmargin=0pt,
  aboveskip=1.0em,
  belowskip=1.0em,
}
\lstset{style=py}

\newcommand{\cheb}[1]{T_{#1}}
\newcommand{\R}{\mathbb{R}}
\newcommand{\C}{\mathbb{C}}
\newcommand{\eps}{\varepsilon_{\mathrm{mach}}}
\newcommand{\code}[1]{\texttt{#1}}

\title{\textbf{ChebPy}\\[0.35em]
       \large Computing with Functions via Chebyshev Approximation in Python}
\author{The ChebPy Developers}
\date{Version 0.10.0}

\begin{document}
\maketitle

\begin{abstract}
\noindent
ChebPy is a Python library for numerical computing with \emph{functions} rather
than numbers. A smooth function is replaced by a polynomial interpolant through
Chebyshev points, accurate to close to machine precision, and thereafter
differentiation, integration, rootfinding, convolution and ordinary arithmetic
are carried out on that surrogate. The result is a system in which
\code{f + g}, \code{f.diff()} and \code{f.roots()} mean what a mathematician
expects them to mean. This paper describes the mathematical foundation
(Section~\ref{sec:approx}), the numerical kernels that implement it
(Section~\ref{sec:algorithms}), the two-hierarchy object model that organises
them (Section~\ref{sec:architecture}), the extensions that carry the idea
beyond smooth functions on bounded intervals --- Fourier technology for
periodic functions, support truncation for infinite intervals and
endpoint-clustering maps for branch singularities (Section~\ref{sec:beyond})
--- and the higher-level applications built on top (Section~\ref{sec:apps}).
ChebPy is a Python descendant of the MATLAB package Chebfun, and follows its
algorithmic choices closely while departing from them where a different design
suits the host language better.
\end{abstract}

\tableofcontents
\bigskip

%==============================================================================
\section{Introduction}
\label{sec:intro}

Numerical software usually asks the user to discretise first and compute
afterwards: choose a grid, evaluate, and then apply a quadrature rule or a
finite-difference stencil to the resulting array. The discretisation is the
user's responsibility, and its accuracy is the user's problem.

ChebPy inverts that arrangement. The object of computation is the function
itself. When a user writes

\begin{lstlisting}
from chebpy import chebfun
import numpy as np

f = chebfun(lambda x: np.sin(10 * x) * np.exp(-x**2), [-3, 3])
f.sum()        # definite integral over [-3, 3]
f.diff()       # derivative, again a chebfun
f.roots()      # every root in [-3, 3]
(f * f).sum()  # integral of the square
\end{lstlisting}

\noindent
the constructor samples \code{f} at Chebyshev points, adaptively increasing the
sample count until the Chebyshev coefficients have decayed to the level of
rounding error, and stores the truncated coefficient vector. Every subsequent
operation is exact on that polynomial, and therefore accurate to roughly
machine precision on the original function. The user chose no grid and no step
size.

This is the design of Chebfun~\cite{chebfun,driscoll2014}, the MATLAB system
developed at Oxford since 2004. ChebPy brings it to Python, built on NumPy, and
is distributed on PyPI under the name \code{chebfun} while being imported as
\code{chebpy}. The library was begun by Mark Richardson and is released under
the BSD-3-Clause licence.

The central claim that makes any of this practical is one of approximation
theory: for functions with even modest smoothness, polynomial interpolation in
Chebyshev points converges fast, and for analytic functions it converges
geometrically. Section~\ref{sec:approx} states this precisely; the rest of the
paper is about turning it into software.

%==============================================================================
\section{Chebyshev approximation}
\label{sec:approx}

\subsection{Chebyshev polynomials and series}

The Chebyshev polynomial of the first kind of degree $k$ is defined on
$[-1,1]$ by
\begin{equation}
  \cheb{k}(x) = \cos\!\bigl(k \arccos x\bigr), \qquad x \in [-1, 1],
  \label{eq:Tdef}
\end{equation}
equivalently by the three-term recurrence
\begin{equation}
  \cheb{0}(x) = 1, \quad \cheb{1}(x) = x, \quad
  \cheb{k+1}(x) = 2x\,\cheb{k}(x) - \cheb{k-1}(x).
  \label{eq:Trec}
\end{equation}
Under the substitution $x = \cos\theta$, \eqref{eq:Tdef} becomes
$\cheb{k}(\cos\theta) = \cos k\theta$: a Chebyshev series in $x$ is a cosine
series in $\theta$. Nearly every algorithm in this paper is, at bottom, an
exploitation of that identity --- it is what allows the fast Fourier transform
to do the work.

A Lipschitz-continuous $f$ on $[-1,1]$ has a unique absolutely and uniformly
convergent Chebyshev expansion
\begin{equation}
  f(x) = \sum_{k=0}^{\infty} a_k \cheb{k}(x),
  \qquad
  a_k = \frac{2 - \delta_{k0}}{\pi} \int_{-1}^{1}
        \frac{f(x)\,\cheb{k}(x)}{\sqrt{1 - x^2}} \, dx .
  \label{eq:series}
\end{equation}

\subsection{Why the coefficients decay}

The speed at which $a_k \to 0$ governs everything. Let $E_\rho$ denote the
\emph{Bernstein ellipse}: the image of the circle $|z| = \rho > 1$ under the
Joukowski map $z \mapsto (z + z^{-1})/2$, an ellipse with foci $\pm 1$ and
semi-axis sum $\rho$.

\begin{theorem}[Geometric convergence; see \cite{trefethen2013}, Ch.~8]
\label{thm:analytic}
If $f$ is analytic in $E_\rho$ and satisfies $|f| \le M$ there, then
\[
  |a_k| \le 2M\rho^{-k},
  \qquad
  \bigl\| f - f_n \bigr\|_\infty \le \frac{2M\rho^{-n}}{\rho - 1},
\]
where $f_n$ is the degree-$n$ truncation of \eqref{eq:series}.
\end{theorem}

\begin{theorem}[Algebraic convergence; see \cite{trefethen2013}, Ch.~7]
\label{thm:diff}
If $f$ has $\nu \ge 1$ derivatives on $[-1,1]$ with $f^{(\nu)}$ of bounded
variation $V$, then $|a_k| = O(k^{-\nu-1})$ and
$\|f - f_n\|_\infty = O(n^{-\nu})$.
\end{theorem}

So an entire function is captured to machine precision by a few dozen
coefficients; a $C^\infty$ function by rather more; a function with a branch
point on $[-1,1]$ by far too many --- which is precisely the case handled by
the splitting and singularity machinery of Sections~\ref{sec:chebfun-container}
and~\ref{sec:singfun}.

\subsection{Interpolation, and why it is nearly as good}

ChebPy does not compute the integrals in \eqref{eq:series}. It interpolates.
The \emph{Chebyshev points of the second kind} are
\begin{equation}
  x_j = \cos\!\left(\frac{(n-1-j)\pi}{n-1}\right), \qquad j = 0, \ldots, n-1,
  \label{eq:chebpts}
\end{equation}
which are the extrema of $\cheb{n-1}$ together with the endpoints $\pm 1$,
ordered left to right so that $x_0 = -1$ and $x_{n-1} = 1$. In ChebPy these are
produced by \code{algorithms.chebpts2(n)}.

Let $p_{n-1}$ be the unique polynomial of degree $\le n-1$ interpolating $f$ at
these points, with discrete Chebyshev coefficients $c_k$. The relationship
between $c_k$ and the true $a_k$ is \emph{aliasing}: on the grid
\eqref{eq:chebpts}, $\cheb{m}$ and $\cheb{2(n-1)\pm m}$ are indistinguishable,
so
\begin{equation}
  c_k = a_k + a_{2(n-1)-k} + a_{2(n-1)+k} + a_{4(n-1)-k} + \cdots .
  \label{eq:aliasing}
\end{equation}
Under the hypotheses of Theorem~\ref{thm:analytic} the aliasing error is of the
same order as the truncation error itself. The practical consequence is the one
that matters: \emph{interpolation is within a small factor of the best
polynomial approximation}, so ChebPy can use function samples --- cheap --- and
lose almost nothing.

Two classical facts complete the picture. Interpolation in these points is
stable: the Lebesgue constant grows only like $\tfrac{2}{\pi}\log n$. And it
underlies Clenshaw--Curtis quadrature, so the integration routine of
Section~\ref{sec:calculus} inherits the same accuracy as the representation.

%==============================================================================
\section{The numerical kernels}
\label{sec:algorithms}

This section describes the algorithms in \code{chebpy.algorithms} and
\code{chebpy.chebtech}, which together constitute the computational core.
Table~\ref{tab:complexity} summarises their costs.

\begin{table}[htbp]
\centering
\begin{tabular}{@{}llc@{}}
\toprule
\textbf{Operation} & \textbf{Method} & \textbf{Cost} \\
\midrule
Values $\to$ coefficients & DCT-I via FFT           & $O(n \log n)$ \\
Coefficients $\to$ values & IDCT-I via FFT          & $O(n \log n)$ \\
Truncation length         & \code{standard\_chop}   & $O(n)$ \\
Construction              & adaptive doubling       & $O(n \log n)$ \\
Evaluation at $m$ points  & Clenshaw / barycentric  & $O(mn)$ \\
Differentiation           & coefficient recurrence  & $O(n)$ \\
Indefinite integration    & coefficient recurrence  & $O(n)$ \\
Definite integration      & Clenshaw--Curtis weights& $O(n)$ \\
Multiplication            & FFT coefficient convolution & $O(n \log n)$ \\
Rootfinding               & colleague eigenvalues + subdivision & $O(n^3)$ per block \\
Chebyshev $\leftrightarrow$ Legendre & three-term recurrence & $O(n^2)$ \\
\bottomrule
\end{tabular}
\caption{Cost of the principal kernels, for a representation of length $n$.}
\label{tab:complexity}
\end{table}

\subsection{The discrete transform}

Converting between the $n$ samples $v_j = f(x_j)$ and the $n$ discrete
coefficients $c_k$ is a discrete cosine transform of type I. ChebPy realises it
with NumPy's complex FFT by even extension. In
\code{algorithms.vals2coeffs2}, the sample vector is mirrored to length
$2n - 2$,
\[
  \tilde v = \bigl(v_{n-1}, v_{n-2}, \ldots, v_0,\;
                   v_1, v_2, \ldots, v_{n-2}\bigr),
\]
an inverse FFT is applied, the first $n$ entries are retained, and the interior
entries $c_1, \ldots, c_{n-2}$ are doubled. The inverse map,
\code{coeffs2vals2}, reverses these steps with a forward FFT. Both routines
detect purely real or purely imaginary input and take the real part
accordingly, so a real function never acquires a spurious imaginary component
through the transform.

\subsection{Deciding where to stop: \texttt{standard\_chop}}
\label{sec:chop}

Given a coefficient vector believed to have converged, how many entries are
genuinely significant? Truncating too early loses accuracy; truncating too late
stores rounding noise and slows every later operation. ChebPy implements the
algorithm of Aurentz and Trefethen~\cite{aurentz2015} in
\code{algorithms.standard\_chop}, which proceeds in three steps.

\begin{enumerate}
  \item \textbf{Envelope.} Form the monotonically non-increasing envelope of
        $|c_k|$ by taking a running maximum from the tail backwards, normalised
        so that its first entry is $1$.
  \item \textbf{Plateau.} Scan for the first index $j$ at which the envelope
        stops decreasing meaningfully --- the point where genuine decay gives
        way to the flat noise floor. The test compares the envelope at $j$ with
        its value at $\lceil 1.25 j + 5 \rceil$ against a tolerance-dependent
        ratio $r = 3\bigl(1 - \log e_1 / \log \varepsilon\bigr)$.
  \item \textbf{Refinement.} Within the plateau, minimise
        $\log_{10}(\text{envelope})$ plus a linear ramp that biases the choice
        leftwards, and cut there.
\end{enumerate}

Vectors shorter than $17$ entries are never chopped: there is not enough
evidence of decay to justify it. The tolerance $\varepsilon^{7/6}$ appears in
step~3 as a floor, deliberately slightly looser than $\varepsilon$ itself.

\subsection{Adaptive construction}
\label{sec:adaptive}

\code{algorithms.adaptive} drives the constructor. It samples on grids of
$n = 2^k + 1$ points for $k = 4, 5, \ldots$ --- that is, $17, 33, 65, \ldots$
--- and at each stage transforms to coefficients and applies
\code{standard\_chop}. If the chop length is strictly less than the vector
length, the coefficients have resolved the function and the loop exits with the
truncated vector. Doubling is used so that the total cost is dominated by the
final, successful transform.

Two guards matter. If the largest sampled value falls below the working
tolerance, the function is declared indistinguishable from zero and represented
by the single coefficient $0$; this avoids the degenerate case in which a
vanishing vertical scale makes every relative test meaningless. And if $k$
reaches \code{prefs.maxpow2} (default $16$, i.e.\ $65\,537$ points) without
resolution, ChebPy emits a \code{UserWarning} naming the class that failed and
returns the unresolved coefficients rather than raising --- a non-converged
answer with a warning being more useful, in interactive work, than no answer.

The working tolerance is $\eps \cdot \max(h, 1)$, where $h$ is a horizontal
scale passed in by the caller, so that a function on a wide interval is not
held to an unattainably tight absolute standard.

\subsection{Evaluation}

Two evaluation schemes coexist, and \code{Chebtech.\_\_call\_\_} selects
between them by keyword (\code{how="clenshaw"} by default).

\paragraph{Clenshaw's algorithm.}
\code{algorithms.clenshaw} sums the series directly through the backward
recurrence associated with \eqref{eq:Trec},
\begin{equation}
  b_k = a_k + 2x\,b_{k+1} - b_{k+2}, \qquad
  b_{n} = b_{n+1} = 0, \qquad
  p(x) = a_0 + x\,b_1 - b_2 .
  \label{eq:clenshaw}
\end{equation}
The implementation unrolls the loop two steps at a time, which halves the
number of vector temporaries NumPy must allocate. The cost is $O(n)$ per
evaluation point with no precomputation.

\paragraph{Barycentric interpolation.}
\code{algorithms.bary} evaluates instead from the samples, using the second
barycentric formula~\cite{berrut2004}
\begin{equation}
  p(x) = \left. \sum_{j=0}^{n-1} \frac{w_j}{x - x_j} f_j \middle/
                \sum_{j=0}^{n-1} \frac{w_j}{x - x_j} \right. ,
  \label{eq:bary}
\end{equation}
which for the points \eqref{eq:chebpts} has the remarkably simple weights
$w_j = (-1)^j$, halved at the two endpoints (\code{algorithms.barywts2}).
Formula \eqref{eq:bary} is numerically stable and $O(n)$ per point. ChebPy
switches between two loop orderings --- over evaluation points when there are
few of them, over nodes when there are many --- and repairs the removable
singularities at $x = x_j$, where \eqref{eq:bary} produces $0/0$, by
substituting the exact nodal value.

\subsection{Calculus}
\label{sec:calculus}

Because differentiation and integration act linearly on \eqref{eq:series} and
$\cheb{k}$ satisfies simple recurrences, all three calculus operations are
$O(n)$ manipulations of the coefficient vector --- no quadrature nodes, no step
size.

\paragraph{Definite integral.}
From $\int_{-1}^{1} \cheb{k}(x)\,dx = 0$ for odd $k$ and $2/(1-k^2)$ for even
$k$,
\begin{equation}
  \int_{-1}^{1} f(x)\,dx
    \;=\; 2a_0 + \sum_{\substack{k \ge 2 \\ k \text{ even}}}
          \frac{2\,a_k}{1 - k^2},
  \label{eq:sum}
\end{equation}
implemented in \code{Chebtech.sum} by zeroing the odd coefficients and
contracting with the weight vector. This is Clenshaw--Curtis quadrature,
expressed in coefficient space.

\paragraph{Indefinite integral.}
\code{Chebtech.cumsum} returns $F$ with $F(-1) = 0$ from
\begin{equation}
  b_k = \frac{a_{k-1} - a_{k+1}}{2k} \;\; (k \ge 2),
  \qquad
  b_1 = a_0 - \tfrac{1}{2} a_2,
  \qquad
  b_0 = \sum_{k \ge 1} (-1)^{k+1} b_k ,
  \label{eq:cumsum}
\end{equation}
the last line being exactly the constant that enforces the boundary condition,
since $\cheb{k}(-1) = (-1)^k$.

\paragraph{Derivative.}
\code{Chebtech.diff} inverts \eqref{eq:cumsum}. Writing $f' = \sum_k z_k
\cheb{k}$,
\begin{equation}
  z_{k-1} = z_{k+1} + 2k\,a_k, \qquad k = n-1, n-2, \ldots, 1,
  \label{eq:diff}
\end{equation}
with $z_0$ subsequently halved. The implementation evaluates the two parity
classes as separate reversed cumulative sums, so the whole recurrence is two
NumPy \code{cumsum} calls rather than a Python loop.

Note that \eqref{eq:diff} multiplies by $k$: differentiation amplifies the
high-order coefficients, and repeated differentiation loses roughly one digit
per application to a well-resolved function. This is intrinsic to the operation,
not to the representation.

\subsection{Multiplication}

The product of two Chebyshev series is a convolution of their coefficients,
which \code{coeffmult} evaluates in $O(n \log n)$. Each coefficient
vector is extended to a symmetric sequence of length $2n-2$ with the leading
entry doubled; the two extensions are multiplied pointwise in Fourier space;
the result is folded back and scaled by $1/4$. Since the product of two
degree-$n$ polynomials has degree $2n$, the operands are first prolonged to a
common length so that the product is represented exactly rather than
implicitly aliased.

\subsection{Rootfinding}
\label{sec:roots}

Finding all roots of a polynomial of degree several hundred is the one place
where a naive approach fails outright, and ChebPy's treatment in
\code{algorithms.rootsunit} is correspondingly careful.

\paragraph{Colleague matrix.}
For a series $f = \sum_{k=0}^{N} a_k \cheb{k}$ with $a_N \ne 0$, the roots are
the eigenvalues of the $N \times N$ \emph{colleague
matrix}~\cite{good1961,boyd2002}
\begin{equation}
  A \;=\;
  \begin{pmatrix}
    0 & 1 & & & \\
    \tfrac{1}{2} & 0 & \tfrac{1}{2} & & \\
    & \ddots & \ddots & \ddots & \\
    & & \tfrac{1}{2} & 0 & \tfrac{1}{2} \\
    & & & \tfrac{1}{2} & 0
  \end{pmatrix}
  \;-\;
  \frac{1}{2 a_N}\,
  e_N \bigl(a_0,\, a_1,\, \ldots,\, a_{N-1}\bigr),
  \label{eq:colleague}
\end{equation}
the Chebyshev-basis analogue of the companion matrix, formed as a symmetric
tridiagonal part plus a rank-one correction in the last row. ChebPy passes $A$
to \code{numpy.linalg.eigvals}.

\paragraph{Subdivision.}
The eigenvalue solve is $O(N^3)$, so it is used only for $N \le 50$. Above that
threshold ChebPy bisects the interval, evaluates the series on each half by
Clenshaw, re-expands each half by \code{vals2coeffs2}, and recurses. The split
point is not $0$ but
\[
  \sigma = -0.004849834917525,
\]
a deliberately irrational-looking constant inherited from Chebfun. Splitting at
the exact midpoint would repeatedly land breakpoints on symmetry axes of even
or odd functions, where roots cluster and the two halves inherit the worst of
the conditioning; an off-centre split avoids that resonance. The recursion
doubles the tolerance at each level, since accumulated map compositions make
the deep sub-intervals progressively less able to sustain the original
accuracy.

\paragraph{Filtering and polishing.}
The eigenvalues of \eqref{eq:colleague} are complex in general. Roots with
$|\operatorname{Im}| \ge h_{\mathrm{tol}}$ are discarded, the survivors are
taken as real, those outside $[-1-h_{\mathrm{tol}}, 1+h_{\mathrm{tol}}]$ are
dropped, and the extreme two are clamped to the interval so that a root at an
endpoint is reported exactly there. \code{algorithms.newtonroots} then applies
Newton polishing, at most \code{prefs.maxiter} $= 10$ iterations, to recover the
digits lost in the eigenvalue solve.

\subsection{Legendre conversion and convolution}
\label{sec:conv}

Convolution is not natural in the Chebyshev basis, but it is in the Legendre
basis, where the Hale--Townsend algorithm~\cite{hale2014} computes
$h = f \star g$ in $O(n^2)$ via a recurrence on the Legendre coefficients.
ChebPy therefore provides \code{cheb2leg} and \code{leg2cheb}, each an $O(n^2)$
three-term recurrence derived from \eqref{eq:Trec} and from
$(k+1)P_{k+1} = (2k+1)xP_k - kP_{k-1}$ respectively, built column by column so
that the recurrence is applied to whole coefficient vectors at once.

\code{chebpy.\_convolution} then dispatches on the shape of the operands. Two
single-piece funs of equal width take the fast Legendre path, whose output is a
piecewise polynomial on $[-2,2]$ split at the origin --- the convolution of two
functions supported on $[-1,1]$ generically has a derivative discontinuity
there, and representing it as two pieces keeps both smooth. Anything else falls
back to building each output sub-interval adaptively by Gauss--Legendre
quadrature.

%==============================================================================
\section{Architecture}
\label{sec:architecture}

\subsection{Two hierarchies meeting at a seam}

The library separates \emph{what the approximation is} from \emph{where it
lives}. Two runtime type hierarchies express that split and meet at
\code{Classicfun}:

\begin{center}
\small
\begin{tabular}{@{}>{\raggedright\arraybackslash}p{0.56\textwidth}
                  >{\raggedright\arraybackslash}p{0.36\textwidth}@{}}
\toprule
\textbf{Hierarchy} & \textbf{Meaning} \\
\midrule
\code{Onefun $\to$ Smoothfun $\to$\newline
      \{Chebtech, Trigtech\}}
  & a function on the reference interval $[-1,1]$ \\[0.5em]
\code{Fun $\to$ Classicfun $\to$\newline
      \{Bndfun, CompactFun, Singfun\}}
  & that function placed on $[a,b]$ \\
\bottomrule
\end{tabular}
\end{center}

\noindent
A \code{Chebtech} knows about Chebyshev coefficients and nothing about
intervals. A \code{Bndfun} owns a \code{Chebtech} plus an affine map
$[-1,1] \to [a,b]$, and translates between them: evaluation composes with the
forward map, \code{diff} divides by the constant Jacobian, \code{sum}
multiplies by it. Because the Jacobian of an affine map is a constant, these
translations are one-line adjustments --- which is exactly why
Section~\ref{sec:singfun}, where the map is \emph{not} affine, has to override
them.

Above both sits \code{Chebfun}, which is not a \code{Fun} at all but a
\emph{container} of \code{Fun} pieces.

\subsection{Import layering}

The package is organised into strict import layers; a lower layer never imports
an upper one at module scope.

\begin{center}
\small
\begin{tabular}{@{}>{\raggedright\arraybackslash}p{0.34\textwidth}
                  >{\raggedright\arraybackslash}p{0.58\textwidth}@{}}
\toprule
\textbf{Layer} & \textbf{Modules} \\
\midrule
high-level API / applications & \code{api}, \code{quasimatrix}, \code{gpr} \\
piecewise container           & \code{chebfun} (+ \code{\_convolution},
                                \code{\_pointwise},
                                \code{\_singular\_construction},
                                \code{\_ufuncs}) \\
\code{Classicfun} pieces      & \code{bndfun}, \code{compactfun},
                                \code{singfun} \\
interval-mapped base          & \code{classicfun} \\
reference-interval technology & \code{chebtech}, \code{trigtech} \\
numerical kernels             & \code{algorithms}, \code{chebyshev} \\
primitives                    & \code{utilities}, \code{plotting},
                                \code{smoothfun} \\
foundations                   & \code{settings}, \code{exceptions},
                                \code{decorators}, \code{maps},
                                \code{fun}, \code{onefun} \\
\bottomrule
\end{tabular}
\end{center}

\noindent
The foundation layer has no intra-package imports at all. \code{Chebfun}'s
public methods are thin wrappers that delegate to the private implementation
modules, which keeps the main class readable at roughly a thousand lines
despite the breadth of the operator surface.

\subsection{The piecewise container}
\label{sec:chebfun-container}

A single polynomial cannot represent $|x|$ or a step, but a \emph{piecewise}
polynomial can. A \code{Chebfun} holds an ordered list of \code{Fun} pieces on
abutting sub-intervals, together with breakpoint data recording the value at
each junction. Operations distribute over the pieces: \code{sum} adds the piece
integrals, \code{roots} concatenates the piece roots and (under
\code{prefs.mergeroots}) merges duplicates arising at shared breakpoints, and a
binary operation between two Chebfuns first refines both onto the union of
their breakpoints so that the operands are piecewise-aligned.

Breakpoints therefore serve two purposes at once: they carry genuine
discontinuities, and they subdivide regions where a single polynomial would
need an impractical degree.

\subsection{The public surface}

\begin{center}
\begin{tabular}{@{}ll@{}}
\toprule
\textbf{Entry point} & \textbf{Purpose} \\
\midrule
\code{chebfun(f, domain, n=None)} & construct from a callable, adaptively or at
                                    fixed length \\
\code{trigfun(f, domain)}         & construct with Fourier technology
                                    (Section~\ref{sec:trig}) \\
\code{equifun(values, domain)}    & construct from equispaced samples \\
\code{pwc(domain, values)}        & piecewise-constant Chebfun \\
\code{chebpts(n, domain)}         & Chebyshev points and barycentric weights \\
\code{gpr(x, y, ...)}             & Gaussian process regression
                                    (Section~\ref{sec:gpr}) \\
\code{Quasimatrix}, \code{polyfit}& continuous linear algebra
                                    (Section~\ref{sec:quasi}) \\
\code{UserPreferences}            & tolerances and algorithmic switches \\
\bottomrule
\end{tabular}
\end{center}

\noindent
Preferences are a context-managed singleton, so a block of code can tighten or
loosen tolerances locally without disturbing global state:

\begin{lstlisting}
from chebpy import UserPreferences

with UserPreferences() as prefs:
    prefs.eps = 1e-10       # looser tolerance, cheaper constructions
    g = chebfun(f, [0, 1])  # built at 1e-10
# previous preferences restored here
\end{lstlisting}

\noindent
The defaults are $\varepsilon = \eps$, \code{maxpow2} $= 16$,
\code{maxiter} $= 10$ Newton steps, and $2001$ plotting samples.

\subsection{Interoperability}

\code{Chebfun} registers NumPy ufunc overloads (\code{chebpy.\_ufuncs}), so
\code{np.sin(f)}, \code{np.exp(f)} and their relatives return Chebfuns
constructed by composition rather than arrays of samples. The composed function
is passed back through the adaptive constructor, so $\sin(f)$ carries its own
resolved degree rather than inheriting $f$'s.

%==============================================================================
\section{Beyond smooth functions on bounded intervals}
\label{sec:beyond}

Chebyshev interpolation resolves functions that are analytic on a
neighbourhood of a bounded real interval. Three extensions widen that domain of
applicability, each by changing the representation rather than the algorithms
built on it.

\subsection{Periodic functions: \texttt{Trigtech}}
\label{sec:trig}

For a smooth \emph{periodic} function, a Fourier series is the natural basis: it
needs roughly $\pi/2$ times fewer degrees of freedom than a Chebyshev series of
comparable accuracy, and it enforces periodicity exactly instead of
approximately. \code{Trigtech} sits beside \code{Chebtech} under
\code{Smoothfun}, exposing the identical \code{Onefun} interface.

Samples are taken at the $n$ equispaced points $x_j = -1 + 2j/n$, and
coefficients are stored in NumPy's native FFT ordering,
\begin{equation}
  c_k = \frac{1}{n} \sum_{j=0}^{n-1} f(x_j)\, e^{-2\pi i j k / n}
      = \bigl(\texttt{numpy.fft.fft}(v)/n\bigr)_k ,
  \label{eq:trigcoeffs}
\end{equation}
with the constant term at index $0$, positive frequencies next and negative
frequencies wrapped to the end. Evaluation at arbitrary $x \in [-1,1]$ uses the
summation formula
\begin{equation}
  f(x) = \sum_k c_k \, e^{i \pi \omega_k (x + 1)},
  \qquad \omega_k = \texttt{fftfreq}(n) \cdot n .
  \label{eq:trigeval}
\end{equation}
Storing in FFT order rather than the human-readable DC-centred order means the
transforms are direct NumPy calls with no shifting; a helper converts to
centred order for display and coefficient plots.

\subsection{Infinite intervals: \texttt{CompactFun}}
\label{sec:compact}

MATLAB Chebfun handles $[a,\infty)$ and $(-\infty,\infty)$ with
\code{@unbndfun}, mapping the infinite interval onto $[-1,1]$ by a rational
change of variables. ChebPy takes a different route, and the departure is
deliberate.

A \code{CompactFun} distinguishes two intervals. The \emph{logical} interval is
what the user asked for and may have infinite endpoints. The \emph{numerical
support} is the finite region outside which the function differs from its
asymptotic limit by less than a tolerance. The constructor discovers the latter
by probing outwards from a finite anchor --- the bounded endpoint for a
semi-infinite interval, or the origin for the doubly-infinite case --- doubling
the radius up to \code{numsupp\_max\_width} $= 10^6$ over at most
\code{numsupp\_max\_probes} $= 30$ probes, and detecting where the tail has
settled. Inside that region an ordinary \code{Chebtech} is stored; outside it
the function is reported as a scalar constant, \code{tail\_left} or
\code{tail\_right} (default $0$).

The advantage is that the stored object is an ordinary Chebyshev
representation, so every kernel of Section~\ref{sec:algorithms} applies
unchanged and no rational map degrades the conditioning. The limitation is
equally clear: the approach suits functions that decay to a constant, and
integrals that diverge are detected and raise
\code{DivergentIntegralError} rather than returning a meaningless finite
number. For a finite logical interval, storage and logical intervals coincide
and a \code{CompactFun} behaves exactly as a \code{Bndfun}.

\subsection{Endpoint singularities: \texttt{Singfun}}
\label{sec:singfun}

A function such as $\sqrt{x}$ on $[0,1]$ or $\sqrt{x(1-x)}$ is analytic on the
open interval but has a branch point at an endpoint. Its Chebyshev coefficients
decay only algebraically (Theorem~\ref{thm:diff}), so the adaptive constructor
runs to its iteration limit and returns a poor answer.

\code{Singfun} removes the singularity by a change of variable. It is a sibling
of \code{Bndfun} whose only structural novelty is that the bijection between
the storage variable $t \in [-1,1]$ and the logical variable $x \in [a,b]$ is
not affine but a \emph{slit-strip map} of Adcock and
Richardson~\cite{adcock2013}, whose derivative vanishes super-algebraically at
the clustered endpoint. For a left-clustered map with parameters $(L, \alpha)$,
\begin{equation}
  s(y) = \tfrac{L}{2}(y - 1), \qquad
  \gamma = \tfrac{\alpha}{\pi}\log\!\bigl(e^{\pi/\alpha} - 1\bigr), \qquad
  u(s) = \tfrac{\alpha}{\pi}
         \log\!\bigl(1 + e^{\pi (s + \gamma)/\alpha}\bigr),
  \label{eq:slit}
\end{equation}
with $x = a + (b-a)\,u(s(y))$; the shift $\gamma$ is chosen so that $u(0) = 1$.
\code{SingleSlitMap} clusters at one endpoint, \code{DoubleSlitMap} at both.

The composition $f \circ m$ is then analytic in a Bernstein ellipse around
$[-1,1]$ and is resolved by ordinary Chebyshev interpolation to spectral
accuracy. Everything that depends only on the map being a bijection ---
evaluation, rootfinding, the binary operators --- is inherited from
\code{Classicfun} unchanged. Only \code{sum}, \code{cumsum} and \code{diff}
need bespoke implementations, because the constant-Jacobian shortcuts of
Section~\ref{sec:architecture} no longer hold.

The map does not quite cover $[a,b]$: the image of $[-1,1]$ omits a small gap
of relative width \code{gap\_unit} at the clustered end, controlled by $L$.
This is the price of the construction, and the parameters are exposed through
\code{MapParams} so a user can trade coverage against conditioning.

%==============================================================================
\section{Applications}
\label{sec:apps}

\subsection{Quasimatrices}
\label{sec:quasi}

A \emph{quasimatrix} is an $\infty \times n$ matrix: $n$ columns, each a
Chebfun on a common domain, with rows indexed by a continuum. The inner product
$\int_a^b f(x) g(x)\, dx$ replaces the discrete dot product, and with that
substitution much of dense linear algebra carries over. ChebPy's
\code{Quasimatrix} implements QR by the Householder triangularisation of
Trefethen~\cite{trefethen2010}, and from it the SVD, least-squares solves and
polynomial fitting via \code{polyfit}.

The continuous QR is the useful primitive: applied to the quasimatrix whose
columns are $1, x, x^2, \ldots$, Gram--Schmidt would break down numerically at
modest degree, whereas Householder triangularisation returns a numerically
orthonormal basis --- the Legendre polynomials, up to scaling --- and does so
stably.

\subsection{Gaussian process regression}
\label{sec:gpr}

\code{chebpy.gpr} implements Gaussian process regression following Rasmussen
and Williams~\cite{rasmussen2006} and the Chebfun \code{gpr} routine. Given
observations $(x_i, y_i)$, it forms the posterior over functions and returns
the posterior mean as a Chebfun, the variance as a Chebfun, and any requested
posterior samples as a Quasimatrix.

Returning Chebfuns rather than sampled arrays is the point. The posterior mean
can immediately be differentiated, integrated or root-found; the credible band
is a pair of Chebfuns whose crossings with a threshold are located by
Section~\ref{sec:roots}'s rootfinder rather than by scanning a grid. A periodic
kernel is available, in which case the posterior is built on Fourier technology
(Section~\ref{sec:trig}).

%==============================================================================
\section{Implementation and validation}
\label{sec:engineering}

ChebPy targets Python 3.11--3.14 and depends only on NumPy and Matplotlib. Its
development infrastructure is synchronised from a shared template, keeping the
repository's own source clearly separated from generated configuration.

The correctness bars are enforced mechanically rather than by convention. The
test suite spans 23 modules covering every layer from the FFT kernels to the
public API, and coverage is held at \textbf{100\%}; genuinely unreachable
defensive branches must carry an explicit \code{pragma} with a stated reason,
which makes each exemption reviewable rather than invisible. Docstring coverage
is likewise held at 100\%, and the package type-checks cleanly under
\code{mypy --strict} with two independent checkers in the continuous
integration gate. Additional workflows run mutation testing, fuzzing,
CodeQL analysis and dependency auditing.

For a library whose value proposition is ``accurate to machine precision'', the
significance of these bars is specific: a silent inaccuracy in a kernel such as
\code{standard\_chop} or the colleague matrix would not crash anything. It
would return plausible numbers that are wrong in the last several digits. Full
branch coverage of the numerical paths is the practical defence.

%==============================================================================
\section{Conclusion}

ChebPy makes functions into first-class objects by pairing one theorem with a
handful of carefully chosen algorithms. The theorem is that Chebyshev
interpolation converges geometrically for analytic functions and is
near-optimal among polynomial approximations. The algorithms --- the FFT-based
transform, the \code{standard\_chop} truncation test, the adaptive doubling
constructor, Clenshaw and barycentric evaluation, the coefficient recurrences
for calculus, and the colleague-matrix rootfinder with off-centre subdivision
--- turn it into a system where \code{f.roots()} returns every root and
\code{f.sum()} returns the integral to fifteen digits, with no grid chosen by
the user.

The architecture separates the reference-interval representation from its
placement on a domain, which is what allows Fourier technology, infinite
intervals and endpoint singularities to be added as new representations rather
than as special cases threaded through the algorithms. The higher-level
applications --- quasimatrices and Gaussian process regression --- then follow
from the same principle applied once more: if functions are objects, then
matrices of functions and distributions over functions are objects too.

%==============================================================================
\begin{thebibliography}{99}

\bibitem{adcock2013}
B.~Adcock and M.~Richardson,
\emph{New exponential variable transform methods for functions with endpoint
singularities},
arXiv:1305.2643, 2013.

\bibitem{aurentz2015}
J.~L. Aurentz and L.~N. Trefethen,
\emph{Chopping a Chebyshev series},
arXiv:1512.01803, 2015.

\bibitem{berrut2004}
J.-P. Berrut and L.~N. Trefethen,
\emph{Barycentric Lagrange interpolation},
SIAM Review, 46(3):501--517, 2004.

\bibitem{boyd2002}
J.~P. Boyd,
\emph{Computing zeros on a real interval through Chebyshev expansion and
polynomial rootfinding},
SIAM Journal on Numerical Analysis, 40(5):1666--1682, 2002.

\bibitem{clenshaw1955}
C.~W. Clenshaw,
\emph{A note on the summation of Chebyshev series},
Mathematics of Computation, 9(51):118--120, 1955.

\bibitem{driscoll2014}
T.~A. Driscoll, N.~Hale and L.~N. Trefethen (eds.),
\emph{Chebfun Guide},
Pafnuty Publications, Oxford, 2014.

\bibitem{good1961}
I.~J. Good,
\emph{The colleague matrix, a Chebyshev analogue of the companion matrix},
Quarterly Journal of Mathematics, 12(1):61--68, 1961.

\bibitem{hale2014}
N.~Hale and A.~Townsend,
\emph{An algorithm for the convolution of Legendre series},
SIAM Journal on Scientific Computing, 36(3):A1207--A1220, 2014.

\bibitem{rasmussen2006}
C.~E. Rasmussen and C.~K.~I. Williams,
\emph{Gaussian Processes for Machine Learning},
MIT Press, 2006.

\bibitem{trefethen2000}
L.~N. Trefethen,
\emph{Spectral Methods in MATLAB},
SIAM, Philadelphia, 2000.

\bibitem{trefethen2010}
L.~N. Trefethen,
\emph{Householder triangularization of a quasimatrix},
IMA Journal of Numerical Analysis, 30(4):887--897, 2010.

\bibitem{trefethen2013}
L.~N. Trefethen,
\emph{Approximation Theory and Approximation Practice},
SIAM, Philadelphia, 2013.

\bibitem{chebfun}
The Chebfun Developers,
\emph{Chebfun --- numerical computing with functions},
\url{http://www.chebfun.org/}.

\end{thebibliography}

\end{document}
